\subsection{Great Orthogonality Theorem}

\begin{theorem}[Great Orthogonality Theorem]
    Als \(\hat{N}\) alle inequivalente irreps van een eindige groep \(G\) zijn, dan geldt:
    \[
    \frac{1}{|G|}\sum_{x \in G} D^{(\alpha)}_{ij}(x) \bar{D}^{(\beta)}_{kl}(x) = \frac{1}{d_\alpha} \delta_{\alpha\beta} \delta_{ik} \delta_{jl}
    \]
    met \(d_\alpha\) de dimensie van de irrep \(\alpha\).

    Dit betekent dat alle irreps samen een \(\sum_{\alpha=1}^{\hat{N}} d_\alpha^2\)-dimensionale orthogonale basis opspannen in een \(|G|\)-dimensionale vectorruimte.
\end{theorem}

\begin{proof}
    Neem een arbitraire \(d_\alpha \times d_\beta\) matrix \(B\) en definieer
    \[
    Q = \sum_{x\in G} D^\alpha(x)\, B\, D^{\beta\dagger}(x).
    \]

    \textbf{Schur's lemma.}
    Voor elke \(y \in G\) geldt:
    \begin{align*}
        D^\alpha(y)\, Q
        &= \sum_x D^\alpha(yx)\, B\, D^{\beta\dagger}(x) \\
        &= \sum_{x'} D^\alpha(x')\, B\, D^{\beta\dagger}(y^{-1}x') \\
        &= \sum_{x'} D^\alpha(x')\, B\, \left(D^\beta(y^{-1})D^\beta(x')\right)^\dagger \\
        &= \sum_{x'} D^\alpha(x')\, B\, D^{\beta\dagger}(x')\, D^{\beta\dagger}(y^{-1}) \\
        &= \sum_{x'} D^\alpha(x')\, B\, D^{\beta\dagger}(x')\, D^{\beta}(y) \\
        &= Q\, D^\beta(y),
    \end{align*}
    waarbij we de substitutie \(x' = yx\) en unitariteit \(D^{\beta\dagger}(y^{-1}) = D^\beta(y)\) hebben gebruikt.

    Dit is precies de voorwaarde van Schur's lemma, dus
    \[
    Q = \delta_{\alpha\beta} \cdot c_\alpha \cdot \mathds{1}_{d_\alpha}
    \]
    voor een constante \(c_\alpha\) (en \(Q=0\) als \(\alpha \neq \beta\), want dan zijn de domein- en beelddimensies verschillend of de representaties inequivalent).

    \medskip
    \textbf{Bepaal \(c_\alpha\) via het spoor (geval \(\alpha = \beta\)).}

    Aan de ene kant, omdat \(Q = c_\alpha \mathds{1}_{d_\alpha}\):
    \[
    \text{Tr}(Q) = c_\alpha\, d_\alpha.
    \]

    Aan de andere kant, met behulp van de \emph{cykliciteit} van het spoor en unitariteit van \(D^\alpha\) (zodat \(D^{\alpha\dagger}(x)D^\alpha(x) = \mathds{1}\)):
    \begin{align*}
        \text{Tr}(Q) 
        &= \sum_{x \in G} \text{Tr}\left(D^\alpha(x)\, B\, D^{\alpha\dagger}(x)\right) \\
        &= \sum_{x \in G} \text{Tr}\left(D^{\alpha\dagger}(x)\, D^\alpha(x)\, B\right) \\
        &= \sum_{x \in G} \text{Tr}(B) \\
        &= |G| \cdot \text{Tr}(B).
    \end{align*}

    Gelijkstellen geeft:
    \[
    c_\alpha = \frac{|G|\, \text{Tr}(B)}{d_\alpha}.
    \]

    Samengevat:
    \[
    Q_{ik} = \sum_{x \in G} \sum_{k,l} D^\alpha_{ik}(x)\, B_{jl}\, \bar{D}^\beta_{jl}(x)
    = \delta_{\alpha\beta}\,\frac{|G|\,\text{Tr}(B)}{d_\alpha}\,\delta_{ik}.
    \]
    Deze relatie geldt voor \emph{elke} matrix \(B\). Kies nu specifiek \(B =  \ket{l}\bra{j}\),
    dit is de matrix met \((\ket{l}\bra{j})_{ik} = \delta_{ij} \delta_{lk}\):
    \begin{align*}
    Q_{ik} &= \delta_{\alpha\beta} \delta_{lj}\frac{|G|}{d_\alpha} \delta_{ik} \\
    &= \sum_{x \in G} D^\alpha(x)B \bar{D}^\beta(x)_{ik} \\
    &= \sum_{x \in G} \sum_{r,s} (D^\alpha_{ir}(x) B_{rs} \bar{D}^\beta(x))_{ks} \\
    \end{align*}
    \[
    \implies
    \frac{1}{|G|}\sum_{x \in G} D^{(\alpha)}_{ij}(x)\, \bar{D}^{(\beta)}_{kl}(x) = \frac{1}{d_\alpha}\,\delta_{\alpha\beta}\,\delta_{ik}\,\delta_{jl}.
    \]

\end{proof}